\section{From a large-prime tail to the exact density}
\label{sec:sieve}

The geometric argument will estimate large prime-power divisors for a
primitive polynomial with positive leading coefficient. We first prove
that this is enough for the full density statement. The finite-prime
conditions must still be imposed on the original polynomial: removing
its content can change the corresponding factors of the Euler product.

\begin{proposition}[Sieve transfer]\label{prop:sieve-transfer}
Let $f\in\Z[x]$ be irreducible over $\Q$, of degree $d\ge2$, and let
$k\ge2$. Suppose that $\rho_f(p^k)<p^k$ for every prime $p$.
Write
\[
 f=s c g,\qquad s\in\{-1,1\},\quad c\in\N,
\]
where $c$ is the positive content of $f$ and $g\in\Z[x]$ is primitive
with positive leading coefficient. Assume that
\begin{equation}
 T_g(X):=\#\{n\in\Z:X<n\le2X,\ p^k\mid g(n)
                 \text{ for some prime }p>X\}=o_{g,k}(X).
 \label{eq:sieve-tail-hypothesis}
\end{equation}
Then
\[
 S_{f,k}(N)=c_{f,k}N+o_{f,k}(N),\qquad
 c_{f,k}=\prod_p\left(1-\frac{\rho_f(p^k)}{p^k}\right)>0.
\]
\end{proposition}

\begin{proof}
\emph{Normalization and local factors.}
Multiplication by $s$ does not affect divisibility. The polynomial $g$
is irreducible over $\Q$, and it inherits local admissibility from $f$.
Neither polynomial has an integer root, because its degree is at least
two. Thus none of the counted values is zero.

For a prime $p\mid c$, put $e=v_p(c)$. Local admissibility forces $e<k$,
and
\begin{equation}
 \rho_f(p^k)=p^e\rho_g(p^{k-e}),\qquad
 1-\frac{\rho_f(p^k)}{p^k}
 =1-\frac{\rho_g(p^{k-e})}{p^{k-e}}.
 \label{eq:sieve-content-factor}
\end{equation}
Indeed, $p^k\mid f(a)$ is equivalent to $p^{k-e}\mid g(a)$, and each
class modulo $p^{k-e}$ has exactly $p^e$ lifts modulo $p^k$.
For $p\nmid c$ the root counts of $f$ and $g$ modulo $p^k$ agree.
Once $X$ exceeds every prime divisor of $c$, the large-prime exceptional
sets for $f$ and $g$ are therefore identical.

Let $\mathcal B$ contain the prime divisors of $c$, the leading coefficient
of $g$, and the discriminant of $g$. This is a finite set. If
$p\notin\mathcal B$, reduction of $g$ has degree $d$ and simple roots.
Each root modulo $p^j$ lifts uniquely modulo $p^{j+1}$: the expansion
\[
 g(a+p^j u)\equiv g(a)+p^j u g'(a)\pmod{p^{j+1}}
\]
determines $u$ modulo $p$, since $g'(a)$ is invertible there. Consequently
\begin{equation}
 \rho_f(p^k)=\rho_g(p^k)=\rho_g(p)\le d
 \qquad(p\notin\mathcal B).
 \label{eq:sieve-root-bound}
\end{equation}

Put $q_p=\rho_f(p^k)/p^k$. Each $q_p<1$ by hypothesis, and
\eqref{eq:sieve-root-bound} gives $\sum_p q_p<\infty$ because $k\ge2$.
For all sufficiently large $p$, $q_p\le1/2$ and
$-2q_p\le\log(1-q_p)\le-q_p$. Thus the Euler product converges to a
strictly positive number. Its finitely many factors at primes in
$\mathcal B$ are retained exactly, including
\eqref{eq:sieve-content-factor}.

\emph{The finite sieve and the intermediate primes.}
Fix an integer $Y\ge2$ exceeding every prime in $\mathcal B$, and set
\[
 Q_Y=\prod_{p\le Y}p^k,\qquad
 c_Y=\prod_{p\le Y}\left(1-\frac{\rho_f(p^k)}{p^k}\right).
\]
Let $A_Y(X)$ count the integers $n\in(X,2X]$ for which
$p^k\nmid f(n)$ for every $p\le Y$. The Chinese remainder theorem gives
exactly $Q_Yc_Y$ permitted classes modulo $Q_Y$. Counting each class in
the interval gives
\begin{equation}
 A_Y(X)=c_YX+O(Q_Y).
 \label{eq:sieve-finite-crt}
\end{equation}
Here $Y$ remains fixed as $X$ tends to infinity.

Write $A(X)$ for the count of all $k$-power-free values in $(X,2X]$.
For $X>Y$, the integers removed by primes $Y<p\le X$ number at most
\begin{align}
 I_Y(X)
 &\le\sum_{Y<p\le X}\rho_f(p^k)
                    \left(\frac{X}{p^k}+1\right)\notag\\
 &\le dX\sum_{m>Y}\frac1{m^k}+d\pi(X)
 \le\frac{d}{k-1}XY^{1-k}+d\pi(X).
 \label{eq:sieve-intermediate}
\end{align}
The elementary bound $\pi(X)\ll X/\log X$ is sufficient here.
For all sufficiently large $X$, the remaining exceptional set is counted
by $T_g(X)$, so
\begin{equation}
 0\le A_Y(X)-A(X)\le I_Y(X)+T_g(X).
 \label{eq:sieve-sandwich}
\end{equation}

We can now identify the density. Divide
\eqref{eq:sieve-finite-crt}--\eqref{eq:sieve-sandwich} by $X$ and first
let $X\to\infty$ with $Y$ fixed. The tail hypothesis gives
\[
 c_Y-\frac{d}{k-1}Y^{1-k}
 \le\liminf_{X\to\infty}\frac{A(X)}X
 \le\limsup_{X\to\infty}\frac{A(X)}X
 \le c_Y.
\]
Then let $Y\to\infty$. Since $c_Y\to c_{f,k}$, this proves
\begin{equation}
 A(X)=c_{f,k}X+o_{f,k}(X).
 \label{eq:sieve-dyadic-density}
\end{equation}

\emph{From dyadic intervals to positive inputs.}
Fix an integer $J\ge1$. Apply \eqref{eq:sieve-dyadic-density} to the
finitely many disjoint intervals
\[
 (N/2^{j+1},N/2^j],\qquad 0\le j<J.
\]
For fixed $J$, their total count is
$c_{f,k}N(1-2^{-J})+o_{f,k,J}(N)$. The interval $[1,N/2^J]$ contains
at most $N/2^J$ integers. Hence
\[
 \begin{split}
 c_{f,k}(1-2^{-J})
 &\le\liminf_{N\to\infty}\frac{S_{f,k}(N)}N\\
 &\le\limsup_{N\to\infty}\frac{S_{f,k}(N)}N
 \le c_{f,k}(1-2^{-J})+2^{-J}.
 \end{split}
\]
Letting $J\to\infty$ proves the proposition.
\end{proof}

\begin{remark}\label{rem:sieve-congruences}
A large-prime bound remains valid after imposing any fixed finite
congruence restrictions on $n$, since the restricted exceptional set is
a subset of the unrestricted one. In the finite sieve those restrictions
are combined with the conditions modulo $p^k$ using their least common
multiple as modulus. No independence of overlapping congruence conditions
is needed. In particular, the exact conditions at content primes can be
retained while using the large-prime bound for the primitive polynomial.
\end{remark}

The same normalization also applies when a density theorem for the
primitive polynomial is already available. It does not require a
separate estimate for values in arithmetic progressions.

\begin{lemma}[Removing sign and content]\label{lem:density-content}
Let $f\in\Z[x]$ be irreducible over $\Q$, of degree at least two, and let
$k\ge2$. Suppose that $\rho_f(p^k)<p^k$ for every prime $p$, and write
$f=s c g$, where $s\in\{-1,1\}$, $c$ is the positive content, and $g$ is
primitive with positive leading coefficient. If
\[
 S_{g,k}(N)=c_{g,k}N+o_{g,k}(N),\qquad
 c_{g,k}=\prod_p\left(1-\frac{\rho_g(p^k)}{p^k}\right),
\]
then $S_{f,k}(N)=c_{f,k}N+o_{f,k}(N)$ with the Euler product $c_{f,k}>0$
of Proposition~\ref{prop:sieve-transfer}.
\end{lemma}

\begin{proof}
The normalization and product-convergence arguments above do not use
the tail hypothesis. Thus both Euler products converge, and their
factors are related by \eqref{eq:sieve-content-factor}.
Fix an integer $Y$ exceeding every prime divisor of $c$. For $h=f,g$,
let $B_{h,Y}(N)$ count the integers $1\le n\le N$ avoiding $p^k\mid h(n)$
for every $p\le Y$, and put
\[
 c_{h,Y}=\prod_{p\le Y}\left(1-\frac{\rho_h(p^k)}{p^k}\right).
\]
The finite Chinese remainder calculation gives
$B_{h,Y}(N)=c_{h,Y}N+O(Q_Y)$, with $Q_Y=\prod_{p\le Y}p^k$.

An input counted by $B_{f,Y}(N)-S_{f,k}(N)$ avoids all the small-prime
conditions for $f$ but has a divisor $p^k$ with $p>Y$.
It also avoids the small-prime conditions for $g$, and $p\nmid c$ implies
$p^k\mid g(n)$. Consequently
\[
 0\le B_{f,Y}(N)-S_{f,k}(N)
 \le B_{g,Y}(N)-S_{g,k}(N).
\]
For fixed $Y$, divide by $N$ and let $N\to\infty$. The assumed density
of $g$ yields
\[
 c_{f,Y}-(c_{g,Y}-c_{g,k})
 \le\liminf_{N\to\infty}\frac{S_{f,k}(N)}N
 \le\limsup_{N\to\infty}\frac{S_{f,k}(N)}N
 \le c_{f,Y}.
\]
Letting $Y\to\infty$ proves the claim.
\end{proof}

\begin{proof}[Proof of Corollary~\ref{cor:separable-products}]
Write
\[
 F_i=c_i\prod_{g\in\mathcal G_i}g,\qquad c_i\in\Z\setminus\{0\},
\]
where $\mathcal G_i$ is a set of distinct primitive irreducible polynomials
with positive leading coefficients. Let $\mathcal G=\bigcup_i\mathcal G_i$,
so shared factors are included only once. The joint local condition implies
local admissibility of each $F_i$, and hence of every $g\in\mathcal G$.
For such a factor of degree $d\le k+2$, let $T_{g,k}(X)$ count the integers
$X<n\le2X$ for which $p^k\mid g(n)$ for some prime $p>X$. We first record
\begin{equation}
 T_{g,k}(X)=o_{g,k}(X).
 \label{eq:factorwise-tail}
\end{equation}

For $d\le k$, omit the finitely many integer roots of $g$ and use
$|g(n)|\ll_g X^d$. If $d<k$, the tail is empty for large $X$. If $d=k$,
only primes $X<p\le C_gX$ can occur, for a fixed $C_g$. Outside finitely
many primes, each has at most $d$ roots modulo $p^k$, and an interval of
length $X<p^k$ contains at most one representative of each root. Thus
$T_{g,k}(X)\ll_g\pi(C_gX)=o_g(X)$.
If $d=k+1$, the unweighted triple estimates of Reuss
\cite[Lemma~3 and its proof, and Section~6]{Reuss2015PowerFree} give
$o_g(X)$ triples with $X<n\le2X$, $a>X^{1-\delta}$, $\mu^2(a)=1$, and
$a^{d-1}b=g(n)$, for a suitable fixed $\delta>0$. For large $X$ the
values $g(n)$ here are positive, so $a=p$ and $b=g(n)/p^k$ include every
input counted by $T_{g,k}(X)$.
If $d=k+2\le8$, use Proposition~\ref{prop:large-prime-tail}.
Finally, if $d=k+2\ge9$, then $k=d-2\ge(3d+1)/4$ and $k/d>3/4$.
Xiao's estimate in \cite[Section~9]{Xiao2017PowerFree} is
$E_B(B^{1-\delta})=o_g(B)$ for some fixed $\delta>0$, where
\[
 E_B(\xi)=\#\{1\le n\le B:\ b^k\mid g(n)
                 \text{ for some squarefree integer }b>\xi\}.
\]
Take $B=2X$: eventually $p>X>(2X)^{1-\delta}$, so this exceptional set
contains the tail. This proves \eqref{eq:factorwise-tail} in every case.

Let $\mathcal B$ contain the prime divisors of all $c_i$, leading
coefficients and discriminants of the $g\in\mathcal G$, and resultants
of distinct pairs in $\mathcal G$. This set is finite: all the
discriminants and pairwise resultants are nonzero. For $p\notin\mathcal B$,
no two distinct factors vanish at the same residue modulo $p$, and their
roots are simple. Consequently
\[
 \begin{aligned}
 p^k\mid F_i(n)&\ \Longleftrightarrow\
 p^k\mid g(n)\text{ for some }g\in\mathcal G_i,\\
 \#\Omega_p&\le\sum_{g\in\mathcal G}\rho_g(p^k)
 =\sum_{g\in\mathcal G}\rho_g(p)\le D,
 \end{aligned}
\]
where $D=\sum_{g\in\mathcal G}\deg g$. Once $X$ exceeds all primes in
$\mathcal B$, the large-prime exceptional set for the family is contained
in the union of the sets counted by the $T_{g,k}(X)$, and hence has size
$o_{\mathbf F,k}(X)$.

For fixed $Y$ beyond $\mathcal B$, the Chinese remainder theorem leaves
exactly
\[
 Q_Y\prod_{p\le Y}\left(1-\frac{\#\Omega_p}{p^k}\right)
 \quad\text{classes modulo }Q_Y=\prod_{p\le Y}p^k.
\]
The estimate \eqref{eq:sieve-intermediate}, with $D$ in place of $d$,
bounds the contribution of $Y<p\le X$. Combining it with the family
tail and taking the two limits as in
\eqref{eq:sieve-finite-crt}--\eqref{eq:sieve-dyadic-density} gives the
stated density on positive inputs. The product is positive because its
local factors are positive and $\sum_p\#\Omega_p/p^k<\infty$.
At primes in $\mathcal B$ the full set $\Omega_p$ is retained: in
particular, valuations contributed by different factors of one $F_i$
are not separated there. Taking $r=1$ gives the first assertion with
$\#\Omega_p=\rho_F(p^k)$. Empty products cover nonzero constants, and
the finitely many integer zeros of nonconstant polynomials have no
effect on the asymptotic.
\end{proof}
